% !TeX encoding = UTF-8
% !TeX program = xelatex
% !TeX spellcheck = en_US

\documentclass[degree=bachelor,print-style=oneside]{sudathesis}
  % 学位 degree:
  %   doctor | master | bachelor | postdoc
  % 学位类型 degree-type:
  %   academic（默认）| professional
  % 语言 language
  %   chinese（默认）| english
  % 字体库 fontset
  %   windows | mac | fandol | ubuntu
  % 建议终版使用 Windows 平台的字体编译
  % 打印设置 print-style
  %    oneside | twoside
  % 章节编码设置 sequence-style
  %   science | liberal 默认science
  % 默认使用双面打印，根据自己需要设定双面还是单面打印，如果是提交电子版论文建议使用单面打印消除空白页


%%% 配置setup.tex 和main。tex后运行，记得自己选择导入生成封面和选择参考文献格式


% 论文基本配置，加载宏包等全局配置
\input{setup}

\begin{document}

% 封面，使用清华原版的封面
\maketitle
% 插入封面，针对苏大本科生毕业论文
%\includepdf[pages={1}]{./figures/cover.pdf}
% 生成封面，针对苏大研究生毕业论文中文版
\titlech
% 生成封面，针对苏大研究生毕业论文英文版
\titleeng
% 学位论文指导小组、公开评阅人和答辩委员会名单
% 本科生不需要
%% !TeX root = ../thuthesis-example.tex

\begin{committee}[name={学位论文指导小组、公开评阅人和答辩委员会名单}]

  \newcolumntype{C}[1]{@{}>{\centering\arraybackslash}p{#1}}

  \section*{指导小组名单}

  \begin{center}
    \begin{tabular}{C{3cm}C{3cm}C{9cm}@{}}
      李XX & 教授     & 清华大学 \\
      王XX & 副教授   & 清华大学 \\
      张XX & 助理教授 & 清华大学 \\
    \end{tabular}
  \end{center}


  \section*{公开评阅人名单}

  \begin{center}
    \begin{tabular}{C{3cm}C{3cm}C{9cm}@{}}
      刘XX & 教授   & 清华大学                    \\
      陈XX & 副教授 & XXXX大学                    \\
      杨XX & 研究员 & 中国XXXX科学院XXXXXXX研究所 \\
    \end{tabular}
  \end{center}


  \section*{答辩委员会名单}

  \begin{center}
    \begin{tabular}{C{2.75cm}C{2.98cm}C{4.63cm}C{4.63cm}@{}}
      主席 & 赵XX                  & 教授                    & 清华大学       \\
      委员 & 刘XX                  & 教授                    & 清华大学       \\
          & \multirow{2}{*}{杨XX} & \multirow{2}{*}{研究员} & 中国XXXX科学院 \\
          &                       &                         & XXXXXXX研究所  \\
          & 黄XX                  & 教授                    & XXXX大学       \\
          & 周XX                  & 副教授                  & XXXX大学       \\
      秘书 & 吴XX                  & 助理研究员              & 清华大学       \\
    \end{tabular}
  \end{center}

\end{committee}



% 也可以导入 Word 版转的 PDF 文件
% \begin{committee}[file=figures/committee.pdf]
% \end{committee}


% 使用授权的说明
%\copyrightpage
% 将签字扫描后授权文件 scan-copyright.pdf 替换原始页面
% \copyrightpage[file=scan-copyright.pdf]

% 插入独创性声明和使用授权说明（苏大本科生）
\includepdf[pages={1,2}]{./figures/statement.pdf}

\frontmatter
%\newpage
\pagestyle{cnfrontmatter}
\input{contents/abstract_cn}
\newpage
\pagestyle{enfrontmatter}
\input{contents/abstract_en}
\newpage
\pagestyle{plain}
% 目录
\tableofcontents

% 插图和附表清单
% 本科生的插图索引和表格索引需要移至正文之后、参考文献前
% \listoffiguresandtables  % 插图和附表清单（仅限研究生）
\listoffigures           % 插图清单
\listoftables            % 附表清单

% 符号对照表
\begin{denotation}[3cm]
\item[HPC] 高性能计算 (High Performance Computing)
\item[cluster] 集群
\item[Itanium] 安腾
\item[SMP] 对称多处理
\item[API] 应用程序编程接口
\item[PI] 聚酰亚胺
\item[MPI] 聚酰亚胺模型化合物，N-苯基邻苯酰亚胺
\item[PBI] 聚苯并咪唑
\item[MPBI] 聚苯并咪唑模型化合物，N-苯基苯并咪唑
\item[PY] 聚吡咙
\item[PMDA-BDA]	均苯四酸二酐与联苯四胺合成的聚吡咙薄膜
\item[$\Delta G$] 活化自由能 (Activation Free Energy)
\item[$\chi$] 传输系数 (Transmission Coefficient)
\item[$E$] 能量
\item[$m$] 质量
\item[$c$] 光速
\item[$P$] 概率
\item[$T$] 时间
\item[$v$] 速度
\item[劝学] 君子曰：学不可以已。青，取之于蓝，而青于蓝；冰，水为之，而寒于水。木
  直中绳。輮以为轮，其曲中规。虽有槁暴，不复挺者，輮使之然也。故木受绳则直，金就
  砺则利，君子博学而日参省乎己，则知明而行无过矣。吾尝终日而思矣，不如须臾之所学
  也；吾尝跂而望矣，不如登高之博见也。登高而招，臂非加长也，而见者远；顺风而呼，
  声非加疾也，而闻者彰。假舆马者，非利足也，而致千里；假舟楫者，非能水也，而绝江
  河，君子生非异也，善假于物也。积土成山，风雨兴焉；积水成渊，蛟龙生焉；积善成德，
  而神明自得，圣心备焉。故不积跬步，无以至千里；不积小流，无以成江海。骐骥一跃，
  不能十步；驽马十驾，功在不舍。锲而舍之，朽木不折；锲而不舍，金石可镂。蚓无爪牙
  之利，筋骨之强，上食埃土，下饮黄泉，用心一也。蟹六跪而二螯，非蛇鳝之穴无可寄托
  者，用心躁也。—— 荀况
\end{denotation}


%\newpage
% 正文部分
\mainmatter
\startcomponent chap01
\product main

\chapter{模板说明}

\section{test section}
\input knuth

\subsection{sub section}
\input knuth

\subsubsection{another section}
\input knuth
\stopcomponent

\chapter{中华人民共和国}
\label{cha:china}

\section{其它例子}
\label{sec:other}

在第~\ref{cha:intro} 章中我们学习了贝叶斯公式~(\ref{equ:chap1:bayes})，这里我们复
习一下：
\begin{equation}
\label{equ:chap2:bayes}
p(y|\mathbf{x}) = \frac{p(\mathbf{x},y)}{p(\mathbf{x})}=
\frac{p(\mathbf{x}|y)p(y)}{p(\mathbf{x})}
\end{equation}

\subsection{绘图}
\label{sec:draw}

本模板不再预先装载任何绘图包（如 \pkg{pstricks，pgf} 等），完全由用户来决定。
个人觉得 \pkg{pgf} 不错，不依赖于 Postscript。此外还有很多针对 \LaTeX{} 的
 GUI 作图工具，如 XFig(jFig), WinFig, Tpx, Ipe, Dia, Inkscape, LaTeXPiX,
jPicEdt, jaxdraw 等等。

\subsection{插图}
\label{sec:graphs}

强烈推荐《\LaTeXe\ 插图指南》！关于子图形的使用细节请参看 \pkg{subcaption} 宏包的说明文档。

\subsubsection{一个图形}
\label{sec:onefig}
一般图形都是处在浮动环境中。之所以称为浮动是指最终排版效果图形的位置不一定与源文
件中的位置对应\footnote{This is not a bug, but a feature of \LaTeX!}，这也是刚使
用 \LaTeX{} 同学可能遇到的问题。如果要强制固定浮动图形的位置，请使用 \pkg{float} 宏包，
它提供了 \texttt{[H]} 参数，比如图~\ref{fig:xfig1}。
\begin{figure}[H] % use float package if you want it here
  \centering
  \includegraphics{thu-whole-logo}
  \caption{利用 Xfig 制图}
  \label{fig:xfig1}
\end{figure}

大学之道，在明明德，在亲民，在止于至善。知止而后有定；定而后能静；静而后能安；安
而后能虑；虑而后能得。物有本末，事有终始。知所先后，则近道矣。古之欲明明德于天
下者，先治其国；欲治其国者，先齐其家；欲齐其家者，先修其身；欲修其身者，先正其心；
欲正其心者，先诚其意；欲诚其意者，先致其知；致知在格物。物格而后知至；知至而后
意诚；意诚而后心正；心正而后身 修；身修而后家齐；家齐而后国治；国治而后天下
平。自天子以至于庶人，壹是皆以修身为本。其本乱而未治者 否矣。其所厚者薄，而其所
薄者厚，未之有也！

\hfill —— 《大学》


\subsubsection{多个图形}
\label{sec:multifig}

如果多个图形相互独立，并不共用一个图形计数器，那么
用 \texttt{minipage} 或者\texttt{parbox} 就可以。否则，请参看
图~\ref{fig:big1-subcaptionbox}，它包含两个小图，分别是图~\ref{fig:subfig1}和
图~\ref{fig:subfig2}。推荐使用\cs{subcaptionbox}，因为可以像
图~\ref{fig:big1-subcaptionbox} 那样对齐子图的标题，也可以使
用\pkg{subcaption}宏包的\cs{subcaption}（放在 minipage中，用法同\cs{caption}）
或是 \pkg{subfigure} 、 \pkg{subtable}环境，像图~\ref{fig:big1-subfigure}，
不要再用 \cs{subfloat}、\cs{subfigure} 和 \cs{subtable}。

\begin{figure}[h]
  \centering%
  \subcaptionbox{第一个小图形\label{fig:subfig1}}[3cm] %标题的长度，超过则会换行，如下一个小图。
    {\includegraphics[height=3cm]{thu-fig-logo}}%
  \hspace{4em}%
  \subcaptionbox{第二个小图形，注意这个图略矮些。如果标题很长的话，它会自动换行\label{fig:subfig2}}
      {\includegraphics[height=2cm]{thu-text-logo}}
  \caption{包含子图形的大图形(subcaptionbox示例)}
  \label{fig:big1-subcaptionbox}
\end{figure}
\begin{figure}[h]
  \centering%
  \begin{subfigure}{3cm}
    \includegraphics[height=3cm]{thu-fig-logo}
    \caption{第一个小图形}
  \end{subfigure}%
  \hspace{4em}%
  \begin{subfigure}{0.5\textwidth}
    \includegraphics[height=2cm]{thu-text-logo}
    \caption{第二个小图形，注意这个图略矮些。subfigure中同一行的子图在顶端对齐。}
  \end{subfigure}
  \caption{包含子图形的大图形(subfigure示例)}
  \label{fig:big1-subfigure}
\end{figure}

古之学者必有师。师者，所以传道受业解惑也。人非生而知之者，孰能无惑？惑而不从师，
其为惑也，终不解矣。生乎吾前，其闻道也固先乎吾，吾从而师之；生乎吾後，其闻道也亦
先乎吾，吾从而师之。吾师道也，夫庸知其年之先後生於吾乎！是故无贵无贱无长无少，道
之所存，师之所存也。

嗟乎！师道之不传也久矣，欲人之无惑也难矣。古之圣人，其出人也远矣，犹且从师而问焉；
今之众人，其下圣人也亦远矣，而耻学於师。是故圣益圣，愚益愚。圣人之所以为圣，愚
人之所以为愚，其皆出於此乎？爱其子，择师而教之，於其身也，则耻师焉，惑焉。彼童子
之师，授之书而习其句读者，非吾所谓传其道、解其惑者也。句读之不知，惑之不解，或师
焉，或不焉，小学而大遗，吾未见其明也。巫医、乐师、百工之人不耻相师，  士大夫之族
曰“师”曰“弟子”之云者，则群聚而笑之。问之，则曰：彼与彼年相若也，道相似也，位
卑则足羞，官盛则近谀。呜呼！师道之不复，可知矣。巫医、乐师、百工之人。吾子不齿，
今其智乃反不能及，其可怪也欤！圣人无常师。孔子师郯子、苌子、师襄、老聃。郯子之徒，
其贤不及孔子。孔子曰：“三人行，必有我师。”是故弟子不必不如师，师不必贤於弟子。
闻道有先後，术业有专攻，如是而已。

如果要把编号的两个图形并排，那么小页就非常有用了：
\begin{figure}
\begin{minipage}{0.48\textwidth}
  \centering
  \includegraphics[height=2cm]{thu-whole-logo}
  \caption{并排第一个图}
  \label{fig:parallel1}
\end{minipage}\hfill
\begin{minipage}{0.48\textwidth}
  \centering
  \includegraphics[height=2cm]{thu-whole-logo}
  \caption{并排第二个图}
  \label{fig:parallel2}
\end{minipage}
\end{figure}

李氏子蟠，年十七，好古文、六艺，经传皆通习之，不拘於时，学於余。余嘉其能行古
道，作师说以贻之。

\hfill —— 韩愈（唐）

% !TeX root = ../thuthesis-example.tex

\chapter{研究内容}

本研究旨在探讨Apple Watch开发过程中的开放式创新应用，具体研究内容包括以下几个方面：

\section{研究问题}

本研究围绕以下几个关键问题展开：
\begin{itemize}
    \item 苹果公司在Apple Watch开发过程中如何利用开放式创新获取外部资源和技术？
    \item Apple Watch开发过程中的开放式创新策略对其市场成功有何贡献？
    \item 在Apple Watch的开发中，苹果公司与哪些外部合作伙伴进行了合作，这些合作如何影响了产品的创新？
\end{itemize}

\section{研究目标}

本研究的主要目标是：
\begin{itemize}
    \item 分析Apple Watch的开发过程，梳理其关键里程碑和创新点。
    \item 探讨苹果公司在Apple Watch开发过程中采用的开放式创新策略，包括技术引进、合作开发和用户参与等。
    \item 评估开放式创新对Apple Watch技术创新、市场表现和用户满意度的影响。
    \item 对比Apple Watch与其他可穿戴设备在开放式创新应用方面的差异。
\end{itemize}

\section{研究方法}

为实现上述研究目标，本研究将采用以下方法和步骤：
\begin{itemize}
    \item \textbf{文献研究}：通过学术数据库（如Scopus、Google Scholar）和行业报告收集与开放式创新和Apple Watch相关的文献，梳理现有研究成果。
    \item \textbf{案例分析}：通过对Apple Watch开发过程中的关键案例进行深入分析，揭示其开放式创新策略。
    \item \textbf{数据收集}：收集与Apple Watch开放式创新相关的具体数据，包括合作公告、专利信息和用户反馈等。
    \item \textbf{数据分析}：采用定性和定量分析方法，评估开放式创新对Apple Watch成功的影响。
\end{itemize}

\section{预期成果}

通过本研究，预期将取得以下成果：
\begin{itemize}
    \item 系统梳理Apple Watch开发过程中的开放式创新策略和具体应用案例。
    \item 评估开放式创新对Apple Watch技术创新和市场成功的具体贡献。
    \item 提出开放式创新在可穿戴设备开发中的最佳实践和建议。
\end{itemize}

\section{可能的创新点}

本研究可能的创新点包括：
\begin{itemize}
    \item 填补Apple Watch开发过程中的开放式创新应用研究的空白。
    \item 提供关于开放式创新在科技产品开发中的系统性分析和实证研究，为其他企业提供参考和借鉴。
\end{itemize}

\section{研究框架}

本研究将按照以下框架进行：
\begin{itemize}
    \item \textbf{第一章：引言}：介绍研究背景、研究问题和研究目的。
    \item \textbf{第二章：国内外研究现状}：综述开放式创新和Apple Watch产品开发的现有研究，指出研究不足和空白。
    \item \textbf{第三章：研究内容}：详细描述研究问题、研究目标、研究方法和预期成果。
    \item \textbf{第四章：Apple Watch开发中的开放式创新应用}：具体分析Apple Watch开发过程中的开放式创新策略和应用案例。
    \item \textbf{第五章：研究结论与建议}：总结研究发现，提出对苹果公司及其他企业的建议。
\end{itemize}

% !TeX root = ../thuthesis-example.tex

\chapter{研究方案}

\section{研究方法}

本研究将采用以下研究方法和步骤，以系统性地分析和探讨Apple Watch开发过程中的开放式创新应用。

\subsection{文献研究}

首先，通过学术数据库（如Scopus、Google Scholar）和行业报告收集与开放式创新和Apple Watch相关的文献。具体步骤包括：
\begin{itemize}
    \item 确定关键词，如“开放式创新”、“Apple Watch”、“产品开发”等。
    \item 筛选近十年内的相关文献，以确保研究的时效性。
    \item 对筛选出的文献进行分类和综述，梳理现有研究成果和理论框架。
\end{itemize}

\subsection{案例分析}

通过对Apple Watch开发过程中的关键案例进行深入分析，揭示其开放式创新策略。具体步骤包括：
\begin{itemize}
    \item 选择具有代表性的案例，如Apple Watch与外部医疗技术公司的合作、与Nike的联合开发等\citep{apple2024,davidson_assessing_2023}。
    \item 分析每个案例中的开放式创新策略及其实施过程。
    \item 总结案例中的成功经验和面临的挑战，为其他企业提供借鉴\citep{chesbrough_beyond_2006, gehani2016corporate}。
\end{itemize}

\subsection{数据收集}

收集与Apple Watch开放式创新相关的具体数据，包括合作公告、专利信息和用户反馈等。具体步骤包括：
\begin{itemize}
    \item 查找和收集苹果公司发布的官方合作公告和新闻稿\citep{apple2024}。
    \item 利用专利数据库，收集与Apple Watch相关的专利信息，以了解其技术创新点\citep{dahlander_how_2010}。
    \item 通过用户评论和反馈，分析用户对Apple Watch功能和创新的评价，了解市场需求和满意度\citep{davidson_assessing_2023}。
\end{itemize}

\subsection{数据分析}

采用定性和定量分析方法，评估开放式创新对Apple Watch成功的影响。具体步骤包括：
\begin{itemize}
    \item 对收集到的案例和数据进行定性分析，总结开放式创新策略的实施效果和经验教训\citep{chesbrough2016}。
    \item 利用统计方法，对用户反馈和市场数据进行定量分析，评估开放式创新对产品技术创新、市场表现和用户满意度的具体影响\citep{gehani2016corporate,apple2024}。
\end{itemize}

\section{研究框架}

本研究将按照以下框架进行：
\begin{itemize}
    \item \textbf{第一章：引言}：介绍研究背景、研究问题和研究目的。
    \item \textbf{第二章：国内外研究现状}：综述开放式创新和Apple Watch产品开发的现有研究，指出研究不足和空白。
    \item \textbf{第三章：研究内容}：详细描述研究问题、研究目标、研究方法和预期成果。
    \item \textbf{第四章：Apple Watch开发中的开放式创新应用}：具体分析Apple Watch开发过程中的开放式创新策略和应用案例。
    \item \textbf{第五章：研究结论与建议}：总结研究发现，提出对苹果公司及其他企业的建议。
\end{itemize}

%\chapter{测试}
% 其他部分
\backmatter
% 参考文献
\bibliography{reference/refs}  % 参考文献使用 BibTeX 编译(本科生)

% 研究生引入参考文献使用此版本
%\begin{kai}
%\bibliography{reference/refs}
%\end{kai}

%\printbibliography       % 参考文献使用 BibLaTeX 编译


%\chapter{附录}
% 附录
%\fancypagestyle{input}{%
%  \fancyhf{}%
%  \renewcommand\footrulewidth{0pt}%
%    \renewcommand\headrulewidth{0bp}%
%    \fancyfoot[C]{
%        \xiaowu
%      -\thepage-
%    }%
%}% 给下面引入的所有 pdf 标页码
%\includepdfset{pagecommand={\thispagestyle{input}}} %pdfpages宏包插入文件无页码的解决
%\setcounter{secnumdepth}{0}
%\includepdf[addtotoc={1,chapter,1,文献综述,cc,2,section,2,测试,aa},pages={1,2,3}]{./figures/statement.pdf}

% 本科生需要将附录放到声明之后，个人简历之前
\appendix
% % !TeX root = ../main.tex

\begin{survey}
\label{cha:survey}

\title{Title of the Survey}
\maketitle

写出至少 5000 外文印刷字符的调研阅读报告或者书面翻译 1-2 篇(不少于 2 万外文印刷符)。

It is impossible to cover in a single chapter every concept of mathematical
programming.\cite{tex} This chapter introduces only the basic concepts and techniques of
mathematical programming such that readers gain an understanding of them
throughout the book\cite{abrahams99tex,salomon1995advanced}.


\section{Single-Objective Programming}
The general form of single-objective programming (SOP) is written
as follows,
\begin{equation*} % 如果附录中的公式不想让它出现在公式索引中，那就请
                             % 用 equation*
\left\{\begin{array}{l}
\max \,\,f(x)\\[0.1 cm]
\mbox{subject to:} \\ [0.1 cm]
\qquad g_j(x)\le 0,\quad j=1,2,\cdots,p
\end{array}\right.
\end{equation*}
which maximizes a real-valued function $f$ of
$x=(x_1,x_2,\cdots,x_n)$ subject to a set of constraints.

\newcommand\Real{\mathbf{R}}
\newtheorem{mpdef}{Definition}[chapter]
\begin{mpdef}
In SOP, we call $x$ a decision vector, and
$x_1,x_2,\cdots,x_n$ decision variables. The function
$f$ is called the objective function. The set
\begin{equation*}
S=\left\{x\in\Real^n\bigm|g_j(x)\le 0,\,j=1,2,\cdots,p\right\}
\end{equation*}
is called the feasible set. An element $x$ in $S$ is called a
feasible solution.
\end{mpdef}

\newtheorem{mpdefop}[mpdef]{Definition}
\begin{mpdefop}
A feasible solution $x^*$ is called the optimal
solution of SOP if and only if
\begin{equation}
f(x^*)\ge f(x)
\end{equation}
for any feasible solution $x$.
\end{mpdefop}

One of the outstanding contributions to mathematical programming was known as
the Kuhn-Tucker conditions\ref{eq:ktc}. In order to introduce them, let us give
some definitions. An inequality constraint $g_j(x)\le 0$ is said to be active at
a point $x^*$ if $g_j(x^*)=0$. A point $x^*$ satisfying $g_j(x^*)\le 0$ is said
to be regular if the gradient vectors $\nabla g_j(x)$ of all active constraints
are linearly independent.

Let $x^*$ be a regular point of the constraints of SOP and assume that all the
functions $f(x)$ and $g_j(x),j=1,2,\cdots,p$ are differentiable. If $x^*$ is a
local optimal solution, then there exist Lagrange multipliers
$\lambda_j,j=1,2,\cdots,p$ such that the following Kuhn-Tucker conditions hold,
\begin{equation}
\label{eq:ktc}
\left\{\begin{array}{l}
    \nabla f(x^*)-\sum\limits_{j=1}^p\lambda_j\nabla g_j(x^*)=0\\[0.3cm]
    \lambda_jg_j(x^*)=0,\quad j=1,2,\cdots,p\\[0.2cm]
    \lambda_j\ge 0,\quad j=1,2,\cdots,p.
\end{array}\right.
\end{equation}
If all the functions $f(x)$ and $g_j(x),j=1,2,\cdots,p$ are convex and
differentiable, and the point $x^*$ satisfies the Kuhn-Tucker conditions
(\ref{eq:ktc}), then it has been proved that the point $x^*$ is a global optimal
solution of SOP.

\subsection{Linear Programming}
\label{sec:lp}

If the functions $f(x),g_j(x),j=1,2,\cdots,p$ are all linear, then SOP is called
a {\em linear programming}.

The feasible set of linear is always convex. A point $x$ is called an extreme
point of convex set $S$ if $x\in S$ and $x$ cannot be expressed as a convex
combination of two points in $S$. It has been shown that the optimal solution to
linear programming corresponds to an extreme point of its feasible set provided
that the feasible set $S$ is bounded. This fact is the basis of the {\em simplex
  algorithm} which was developed by Dantzig as a very efficient method for
solving linear programming.
\begin{table}[ht]
\centering
  \centering
  \caption*{Table~1\hskip1em This is an example for manually numbered table, which
    would not appear in the list of tables}
  \label{tab:badtabular2}
  \begin{tabular}[c]{|m{1.5cm}|c|c|c|c|c|c|}\hline
    \multicolumn{2}{|c|}{Network Topology} & \# of nodes &
    \multicolumn{3}{c|}{\# of clients} & Server \\\hline
    GT-ITM & Waxman Transit-Stub & 600 &
    \multirow{2}{2em}{2\%}&
    \multirow{2}{2em}{10\%}&
    \multirow{2}{2em}{50\%}&
    \multirow{2}{1.2in}{Max. Connectivity}\\\cline{1-3}
    \multicolumn{2}{|c|}{Inet-2.1} & 6000 & & & &\\\hline
    \multirow{2}{1.5cm}{Xue} & Rui  & Ni &\multicolumn{4}{c|}{\multirow{2}*{\thuthesis}}\\\cline{2-3}
    & \multicolumn{2}{c|}{ABCDEF} &\multicolumn{4}{c|}{} \\\hline
\end{tabular}
\end{table}

Roughly speaking, the simplex algorithm examines only the extreme points of the
feasible set, rather than all feasible points. At first, the simplex algorithm
selects an extreme point as the initial point. The successive extreme point is
selected so as to improve the objective function value. The procedure is
repeated until no improvement in objective function value can be made. The last
extreme point is the optimal solution.

\subsection{Nonlinear Programming}

If at least one of the functions $f(x),g_j(x),j=1,2,\cdots,p$ is nonlinear, then
SOP is called a {\em nonlinear programming}.

A large number of classical optimization methods have been developed to treat
special-structural nonlinear programming based on the mathematical theory
concerned with analyzing the structure of problems.
\begin{figure}[h]
  \centering
  \includegraphics{thu-lib-logo.pdf}
  \caption*{Figure~1\quad This is an example for manually numbered figure,
    which would not appear in the list of figures}
  \label{tab:badfigure2}
\end{figure}

Now we consider a nonlinear programming which is confronted solely with
maximizing a real-valued function with domain $\Real^n$.  Whether derivatives are
available or not, the usual strategy is first to select a point in $\Real^n$ which
is thought to be the most likely place where the maximum exists. If there is no
information available on which to base such a selection, a point is chosen at
random. From this first point an attempt is made to construct a sequence of
points, each of which yields an improved objective function value over its
predecessor. The next point to be added to the sequence is chosen by analyzing
the behavior of the function at the previous points. This construction continues
until some termination criterion is met. Methods based upon this strategy are
called {\em ascent methods}, which can be classified as {\em direct methods},
{\em gradient methods}, and {\em Hessian methods} according to the information
about the behavior of objective function $f$. Direct methods require only that
the function can be evaluated at each point. Gradient methods require the
evaluation of first derivatives of $f$. Hessian methods require the evaluation
of second derivatives. In fact, there is no superior method for all
problems. The efficiency of a method is very much dependent upon the objective
function.

\subsection{Integer Programming}

{\em Integer programming} is a special mathematical programming in which all of
the variables are assumed to be only integer values. When there are not only
integer variables but also conventional continuous variables, we call it {\em
  mixed integer programming}. If all the variables are assumed either 0 or 1,
then the problem is termed a {\em zero-one programming}. Although integer
programming can be solved by an {\em exhaustive enumeration} theoretically, it
is impractical to solve realistically sized integer programming problems. The
most successful algorithm so far found to solve integer programming is called
the {\em branch-and-bound enumeration} developed by Balas (1965) and Dakin
(1965). The other technique to integer programming is the {\em cutting plane
  method} developed by Gomory (1959).

\hfill\textit{Uncertain Programming\/}\quad(\textsl{BaoDing Liu, 2006.2})

\bibliographystyle{plainnat}
\bibliography{ref/refs,ref/appendix}

\end{survey}
       % 本科生：外文资料的调研阅读报告
% % !TeX root = ../thuthesis-example.tex

\begin{translation}
\label{cha:translation}

\title{书面翻译题目}
\maketitle

\tableofcontents


本科生的外文资料书面翻译。


\section{图表示例}

\subsection{图}

附录中的图片示例（图~\ref{fig:appendix-translation-figure}）。

\begin{figure}
  \centering
  \includegraphics[width=0.6\linewidth]{example-image-a.pdf}
  \caption{附录中的图片示例}
  \label{fig:appendix-translation-figure}
\end{figure}


\subsection{表格}

附录中的表格示例（表~\ref{tab:appendix-translation-table}）。

\begin{table}
  \centering
  \caption{附录中的表格示例}
  \begin{tabular}{ll}
    \toprule
    文件名          & 描述                         \\
    \midrule
    thuthesis.dtx   & 模板的源文件，包括文档和注释 \\
    thuthesis.cls   & 模板文件                     \\
    thuthesis-*.bst & BibTeX 参考文献表样式文件    \\
    thuthesis-*.bbx & BibLaTeX 参考文献表样式文件  \\
    thuthesis-*.cbx & BibLaTeX 引用样式文件        \\
    \bottomrule
  \end{tabular}
  \label{tab:appendix-translation-table}
\end{table}


\section{数学公式}

附录中的数学公式示例（公式\eqref{eq:appendix-translation-equation}）。
\begin{equation}
  \frac{1}{2 \uppi \symup{i}} \int_\gamma f = \sum_{k=1}^m n(\gamma; a_k) \mathscr{R}(f; a_k)
  \label{eq:appendix-translation-equation}
\end{equation}


\section{文献引用}

附录\cite{dupont1974bone}中的参考文献引用\cite{merkt1995rotational}示例
\cite{dupont1974bone,merkt1995rotational}。


\appendix

\section{附录}

附录的内容。


% 书面翻译的参考文献
% 默认使用正文的参考文献样式；
% 如果使用 BibTeX，可以切换为其他兼容 natbib 的 BibTeX 样式。
\bibliographystyle{unsrtnat}
% \bibliographystyle{IEEEtranN}

% 默认使用正文的参考文献 .bib 数据库；
% 如果使用 BibTeX，可以改为指定数据库，如 \bibliography{ref/refs}。
\printbibliography

% 书面翻译对应的原文索引
\begin{translation-index}
  \nocite{mellinger1996laser}
  \nocite{bixon1996dynamics}
  \nocite{carlson1981two}
  \bibliographystyle{unsrtnat}
  \printbibliography
\end{translation-index}

\end{translation}
  % 本科生：外文资料的书面翻译
% TODO！本科生启用appendix请勿启用参考文献section，会出现报错，硕博请自行评估是否需要来启用
% !TeX root = ../thuthesis-example.tex

\chapter{补充内容}

附录是与论文内容密切相关、但编入正文又影响整篇论文编排的条理和逻辑性的资料，例如某些重要的数据表格、计算程序、统计表等，是论文主体的补充内容，可根据需要设置。

附录中的图、表、数学表达式、参考文献等另行编序号，与正文分开，一律用阿拉伯数字编码，
但在数码前冠以附录的序号，例如“图~\ref{fig:appendix-figure}”，
“表~\ref{tab:appendix-table}”，“式\eqref{eq:appendix-equation}”等。


\section{插图}

% 附录中的插图示例（图~\ref{fig:appendix-figure}）。

\begin{figure}
  \centering
  \includegraphics[width=0.6\linewidth]{example-image-a.pdf}
  \caption{附录中的图片示例}
  \label{fig:appendix-figure}
\end{figure}


\section{表格}

% 附录中的表格示例（表~\ref{tab:appendix-table}）。

\begin{table}
  \centering
  \caption{附录中的表格示例}
  \begin{tabular}{ll}
    \toprule
    文件名          & 描述                         \\
    \midrule
    thuthesis.dtx   & 模板的源文件，包括文档和注释 \\
    thuthesis.cls   & 模板文件                     \\
    thuthesis-*.bst & BibTeX 参考文献表样式文件    \\
    thuthesis-*.bbx & BibLaTeX 参考文献表样式文件  \\
    thuthesis-*.cbx & BibLaTeX 引用样式文件        \\
    \bottomrule
  \end{tabular}
  \label{tab:appendix-table}
\end{table}


\section{数学表达式}

% 附录中的数学表达式示例（式\eqref{eq:appendix-equation}）。
\begin{equation}
  \frac{1}{2 \uppi \symup{i}} \int_\gamma f = \sum_{k=1}^m n(\gamma; a_k) \mathscr{R}(f; a_k)
  \label{eq:appendix-equation}
\end{equation}


\section{文献引用}

附录\cite{dupont1974bone}中的参考文献引用\cite{zhengkaiqing1987}示例
\cite{dupont1974bone,zhengkaiqing1987}。

\printbibliography


% 致谢
% !TeX root = ../thuthesis-example.tex

\begin{acknowledgements}
  衷心感谢导师×××教授和物理系××副教授对本人的精心指导。他们的言传身教将使我终生受益。

  在美国麻省理工学院化学系进行九个月的合作研究期间，承蒙 Robert Field 教授热心指导与帮助，不胜感激。

  感谢×××××实验室主任×××教授，以及实验室全体老师和同窗们学的热情帮助和支持！

  本课题承蒙国家自然科学基金资助，特此致谢。
\end{acknowledgements}


% 声明
\statement
% 将签字扫描后的声明文件 scan-statement.pdf 替换原始页面
% \statement[file=scan-statement.pdf]
% 本科生编译生成的声明页默认不加页脚，插入扫描版时再补上；
% 研究生编译生成时有页眉页脚，插入扫描版时不再重复。
% 也可以手动控制是否加页眉页脚
% \statement[page-style=empty]
% \statement[file=scan-statement.pdf, page-style=plain]

% 个人简历、在学期间完成的相关学术成果
% 本科生可以附个人简历，也可以不附个人简历
\begin{resume}

  \resumeitem{���˼���}

  1978~��~4~��~30~�ճ����ڻ���ɽˮ������
  1996~��~9~�¿��뻨��ɽ��ѧ����רҵ��
  2000~��~6~�±��Ʊ�ҵ�������ѧѧʿѧλ��
  ͬ��~9~�����Ա����廪��ѧӦ��ħ��ϵ������ʿ����

  \resumeitem{Ŀǰ����ʽ����������}

  \begin{enumerate}[{$[$}1{$]$}]
  \item Harry Potter, Bajie Zhu, Sanzang Tang, Coupling
    Computation of the BEM and FDM in 3D Spirit Extraction,
    Interplanetarian Conference on Super-Sudden Motion, 2001,
    p.716 - 719(SCI).

  \item
    ���ﲨ�أ����˽䣬�����أ���ά�����캢���߽�Ԫ��ȡ��һ��Ԥ����
    ����������ɽ�������¿���2002, Vol38(16), p.207-209.

  \item Harry Potter, Bajie Zhu, Sanzang Tang, A weighted
    average formula for efficient inductance and resistance
    extraction, International Conference on God 2003, Vol.2,
    p996 - 999.

  \item ����գ����˽䣬���ﲨ�أ�ţħ������������
    ��ά�߽�Ԫ�𹿰�������ȡ�е���ЧԤ������������˿����ѧѧ����2004,
    Vol.44(1), p.45-49(EI).
  \end{enumerate}


  \resumeitem{Ŀǰ�ѱ�¼������}

  \begin{enumerate}[{$[$}1{$]$}]
  \item Wukong Sun, Harry Potter, Sanzang Tang,
    Fast 6-D Impedance Extraction of Spirit,
    International Conference on Communication, Circuit and Systems, 2005, ��¼��.

  \item
    ����գ����ﲨ�أ������أ���������������ζ���Ŀ�����ά��
    ������ȡ�㷨��������ѧ������¼��(EI).
  \end{enumerate}

  \resumeitem{Ŀǰ��������}
  \begin{enumerate}[{$[$}1{$]$}]
  \item Harry Potter, A Fast Algorithm for 3-D Inductance
    Extraction Based on Investigation of Open-Magic
    Current, International Conference on Simulation of
    Alchemy Processes and Devices, 2005.
  \end{enumerate}

\end{resume}


% 指导教师/指导小组评语
% 本科生不需要
% !TeX root = ../thuthesis-example.tex

\chapter{指导教师学术评语}

论文提出了……


% 答辩委员会决议书
% 本科生不需要
% !TeX root = ../thuthesis-example.tex

\begin{resolution}

  论文提出了……

  论文取得的主要创新性成果包括：

  1. ……

  2. ……

  3. ……

  论文工作表明作者在×××××具有×××××知识，具有××××能力，论文××××，答辩××××。

  答辩委员会表决，（×票/一致）同意通过论文答辩，并建议授予×××（姓名）×××（门类）学博士/硕士学位。

\end{resolution}


% 本科生的综合论文训练记录表（扫描版）
% \record{file=scan-record.pdf}

\end{document}
